\scanpage{19}
\mspage{16}
% archive sheet number: 18
Si c'est vrai, et pas connu de plus, ce serait quand $\unsure{=|\simeq}$ scandaleux !

En procédant comme pour les morphismes de chgt de base pour les $Rf_*$, on trouve en tous cas des morphismes canoniques, pour $f:I\longrightarrow J$
\[
u_J f_*^{\mathbb D}\longrightarrow f_*^{\mathbb D'}u_I
\]
et de $=$
\[
f_!^{\mathbb D'}u_I\longrightarrow u_J f_!^{\mathbb D}\,,
\]
et la question est de savoir si c'est des iso.
Questions analogues pour les $f^!$, pour imm. \unsure{fermées|fermée}, (\deleted{\unsure{et}} pour $\unsure{f^?|f_!}$, pour $f$ imm. \unsure{ouverte|ouvertes}). Comme la notion de morphisme de dérivateurs devrait quand $=$ ``coller'' à la notion bien connue de morphisme \inserted{(triangulé)} de cat. triangulées, donc être ``exact'' dans un sens convenable, elle montre bien $*$ que l'on doit exiger la commutation des \deleted{\unsure{fonct}} morphismes de dérivateurs aux $f_*$ et aux $f_!$.
% slanted
\margin{Non ! Le cas des cat. triangulées est ``misleading'', car, pour celles-ci, ``exactitude à g. et à dr.'' \unsure{reviennent} au même.}
% slanted
\margin{\uline{NB} On conçoit qu'un foncteur entre dérivateurs puisse être exact à g. ou exact à dr., ou les deux !}

\unsure{J|Je} devrais demander à Thomason s'il a regardé, pour un morphisme
\[
\mathcal M\longrightarrow\mathcal M'
\]
entre cat. de modèles, qui \deleted{soit} \deleted{\unsure{exact}} disons, \uline{exact à gauche}, et ``compatible avec les \inserted{équivalences,} fibrations (et cofibrations'' (???)).
\margin{Notation $\mathbb{DER}$ pour la 2-catégorie des \uline{dérivateurs}}

\figure{\begin{tikzpicture}\draw[line width=.6pt] (0,.20) -- (.16,0);\end{tikzpicture}}
